\subsection{49. The Outer Mucus Layer Hosts a Distinct Intestinal Microbial Niche}
\textbf{PMID:} \texttt{26392213}\quad\textbf{Authors:} Li et al.\ (2015)\quad\textbf{Journal:} Nature Communications 6, 8292\quad\textbf{Domain:} Microbiome / 16S Amplicon\\
\scorebadge{Coverage}{6}\quad\scorebadge{Agreement}{6}\quad\textbf{Total:} 12/20\par\smallskip
\noindent\textbf{Coverage rationale.} 16S amplicon data downloaded from Figshare (3 FASTQ + 3 mapping files; not SRA as implied). 101 SPF samples + 60 sDMDMm2 samples processed. PCoA ordination and PERMANOVA/ANOSIM computed. Critical gap: Bray-Curtis used instead of weighted UniFrac (requires phylogenetic tree + Greengenes taxonomy); truncation-based OTU clustering instead of UCLUST 97\%.\par
\noindent\textbf{Agreement rationale.} PERMANOVA detects significant compartment effect (SPF: $p=0.001$, $R^2=3.0\%$; sDMDMm2: $p=0.003$, $R^2=6.8\%$). However, effect sizes are very small and ANOSIM R values near zero ($\sim$0.05--0.06). Visual PCoA shows extensive overlap, not distinct clustering. The word ``distinct'' in the title overstates the signal.\par\smallskip
\noindent\textbf{Verdict: PARTIAL.} Statistical significance replicates but effect sizes are small; full replication requires weighted UniFrac and proper OTU clustering.

